\documentclass[12pt,a4paper]{article}
\usepackage[margin=1in]{geometry}
\usepackage[T1]{fontenc}
\usepackage{setspace}
\usepackage{amsmath,amssymb}
\usepackage{booktabs}
\usepackage{longtable}
\usepackage{xcolor}
\usepackage{hyperref}
\usepackage{fancyhdr}
\usepackage{listings}

\lstset{basicstyle=\ttfamily\scriptsize,breaklines=true}

\setstretch{1.08}
\setlength{\parskip}{0.45em}
\setlength{\parindent}{0pt}
\setlength{\headheight}{15pt}
\hypersetup{colorlinks=true,linkcolor=blue,citecolor=blue,urlcolor=blue}

\pagestyle{fancy}
\fancyhf{}
\lhead{CST428 - Blockchain Technologies}
\rhead{Module 5 Master Notes}
\cfoot{\thepage}

\newcommand{\examtip}[1]{%
\vspace{0.35em}
\noindent\fcolorbox{black}{gray!12}{\parbox{0.965\linewidth}{\textbf{Exam Tip:} #1}}%
\vspace{0.35em}
}

\begin{document}

\section{Module 5 Topic Map From Previous Questions}

From May 2024, June 2023, and October 2023 papers, Module 5 repeatedly asks:

\begin{itemize}
\item Ethereum transactions: meaning, types, and processing flow.
\item EVM role and Ethereum account types.
\item Ethereum blockchain elements and world state concept.
\item Relationship between transactions, transaction trie, and block header.
\item Gas, gas limit, gas price, and transaction cost impact.
\item Solidity basics: datatypes, control structures, and libraries.
\item Solidity coding questions: bank contract, voting contract, string storage contract.
\item Compare Bitcoin and Ethereum.
\end{itemize}

\section{Core Theory (Explained Clearly)}

\subsection{Ethereum in One Line}
Ethereum is a programmable blockchain where smart contracts execute on a distributed runtime called EVM.

\subsection{EVM (Ethereum Virtual Machine) Made Simple}
Every Ethereum node runs EVM to execute smart contract bytecode in the same way.

Why EVM matters:
\begin{itemize}
\item gives deterministic execution across nodes,
\item updates blockchain state after valid transactions,
\item isolates contract execution with gas metering.
\end{itemize}

\subsection{Ethereum Accounts: Two Types}
\begin{itemize}
\item \textbf{EOA (Externally Owned Account):} controlled by private key, can initiate transactions.
\item \textbf{Contract Account:} controlled by deployed code, executes when called.
\end{itemize}

Quick difference:
\begin{itemize}
\item EOAs have no contract code.
\item Contract accounts have code and persistent storage.
\end{itemize}

\subsection{Ethereum Transactions and Their Types}
A transaction is a signed instruction sent from an EOA to change Ethereum state.

Two major types asked in exams:
\begin{itemize}
\item \textbf{Value transfer transaction:} sends ETH from one account to another.
\item \textbf{Contract interaction/creation transaction:} calls contract function or deploys new contract.
\end{itemize}

Common transaction fields:
\begin{itemize}
\item nonce,
\item to address,
\item value,
\item data payload,
\item gas limit,
\item fee parameters and signature.
\end{itemize}

\subsection{Gas Concept and Formula}
Gas is the unit for computational effort in EVM.

Main formulas:
\[
\text{Transaction Fee} = \text{Gas Used} \times \text{Effective Gas Price}
\]

In EIP-1559 style networks:
\[
\text{Effective Gas Price} = \text{Base Fee} + \text{Priority Fee}
\]

Meaning:
\begin{itemize}
\item More complex function calls consume more gas.
\item If gas limit is too low, transaction fails (reverts) but fee is still spent for computation done.
\end{itemize}

\examtip{In gas answers, always write one numerical-style formula and then explain user impact in plain words.}

\subsection{Ethereum Elements and World State}
Important elements of Ethereum blockchain:
\begin{itemize}
\item block header,
\item transaction list,
\item receipts,
\item state root,
\item transaction root,
\item receipts root,
\item logs and events.
\end{itemize}

World state means a snapshot of all accounts (balances, nonces, contract storage roots, code hashes) at a specific block height.

\begin{center}
\framebox[0.95\linewidth][c]{
\begin{minipage}{0.90\linewidth}
\vspace{0.1cm}
\centerline{\textbf{Ethereum World State (Concept)}}
\centerline{Block N executes transactions}
\centerline{$\Downarrow$}
\centerline{State transitions account-by-account}
\centerline{$\Downarrow$}
\centerline{New state root stored in Block N header}
\centerline{$\Downarrow$}
\centerline{This root defines ``world state after Block N''}
\vspace{0.1cm}
\end{minipage}}
\end{center}

\subsection{Transaction Trie and Block Header Relationship}
Ethereum groups transactions into a Merkle-Patricia Trie.
Its root hash (transaction root) is saved in the block header.

So if one transaction changes, trie root changes, and block header hash changes.
This gives tamper evidence.

\begin{center}
\framebox[0.95\linewidth][c]{
\begin{minipage}{0.90\linewidth}
\vspace{0.1cm}
\centerline{\textbf{Transactions $\rightarrow$ Trie $\rightarrow$ Block Header}}
\centerline{Transaction List}
\centerline{$\Downarrow$}
\centerline{Transaction Trie built from list}
\centerline{$\Downarrow$}
\centerline{Transaction Root inserted into Header}
\centerline{$\Downarrow$}
\centerline{Header hash secures whole block integrity}
\vspace{0.1cm}
\end{minipage}}
\end{center}

\subsection{Solidity Essentials}

\subsubsection*{Datatypes}
\begin{itemize}
\item value types: \texttt{bool}, \texttt{uint}, \texttt{int}, \texttt{address}, \texttt{bytes32}.
\item reference types: arrays, structs, strings.
\item mapping type: key-value storage, for example \texttt{mapping(address => uint)}.
\end{itemize}

\subsubsection*{Control Structures}
\begin{itemize}
\item conditionals: \texttt{if}, \texttt{else}, \texttt{else if}.
\item loops: \texttt{for}, \texttt{while}, \texttt{do-while} (used carefully due to gas risk).
\item flow controls: \texttt{break}, \texttt{continue}, \texttt{return}.
\item error controls: \texttt{require}, \texttt{revert}, \texttt{assert}.
\end{itemize}

\subsubsection*{Libraries in Solidity}
Libraries are reusable code modules deployed once and reused by many contracts.

Benefits:
\begin{itemize}
\item avoids duplicate code,
\item improves maintainability,
\item can reduce deployment size for consuming contracts.
\end{itemize}

\section{Solved Question Papers}

\subsection{Part A: 3-Mark Questions}

\subsubsection*{Q1. Illustrate the role of EVM in Ethereum network.}
\begin{itemize}
\item EVM is the runtime that executes smart contract bytecode on all nodes.
\item It ensures deterministic execution so all honest nodes get same result.
\item It applies state transitions after valid transactions are executed.
\item It isolates code execution and enforces gas-based metering.
\item It is the core reason Ethereum can support programmable contracts.
\item Without EVM, Ethereum would behave like a simple value-transfer chain.
\end{itemize}

\subsubsection*{Q2. What are the two account types in Ethereum? Differentiate them.}
\begin{itemize}
\item Ethereum has EOAs and Contract Accounts.
\item EOAs are controlled by private keys and can start transactions.
\item Contract accounts are controlled by code and cannot self-initiate.
\item EOAs usually store nonce and balance only.
\item Contract accounts store nonce, balance, code, and storage.
\item EOAs call contracts; contracts execute logic when triggered.
\end{itemize}

\subsubsection*{Q3. Explain different datatypes in Solidity.}
\begin{itemize}
\item Value types include \texttt{bool}, \texttt{uint}, \texttt{int}, \texttt{address}, and fixed bytes.
\item Reference types include arrays, structs, and strings.
\item Mapping types store key-value pairs efficiently.
\item Data location matters for references: \texttt{storage}, \texttt{memory}, \texttt{calldata}.
\item Correct datatype choice affects both safety and gas cost.
\item Most exam answers score better when examples are added.
\end{itemize}

\subsubsection*{Q4. Explain how libraries are deployed in Solidity with example idea.}
\begin{itemize}
\item Library code is compiled and deployed as a separate bytecode unit.
\item Contracts can link to the library during deployment.
\item Internal library functions may be inlined by compiler.
\item External library calls happen by delegate-style code execution.
\item Libraries are useful for reusable math and utility operations.
\item Example idea: a \texttt{MathLib} with safe percentage or average functions.
\end{itemize}

\subsubsection*{Q5. Show relation between transactions, transaction trie, and block header.}
\begin{itemize}
\item Block transactions are inserted into a trie structure.
\item Trie root hash summarizes full transaction list.
\item That root hash is written inside block header.
\item Any transaction change alters trie root immediately.
\item Changed trie root changes header hash and block identity.
\item Therefore header indirectly protects all transactions in block.
\end{itemize}

\subsubsection*{Q6. List various components/elements in Ethereum blockchain.}
\begin{itemize}
\item blocks and block headers,
\item transactions and receipts,
\item state trie and state root,
\item transaction trie and transaction root,
\item receipts trie and receipts root,
\item logs/events, accounts, and contract bytecode/storage.
\end{itemize}

\subsection{Part B: 14-Mark Questions}

\subsubsection*{Q1. What is a transaction in Ethereum? Explain two types of transactions in Ethereum.}
\begin{itemize}
\item Ethereum transaction is a signed state-change request sent by an EOA.
\item It is broadcast to the network and validated before inclusion in a block.
\item A transaction carries nonce, destination, value, data, gas limit, fee fields, and signature.
\item Nonce ensures ordering and prevents replay from same account.
\item Value transfer transaction sends ETH between accounts.
\item Contract call transaction sends data to execute function in deployed contract.
\item Contract creation transaction has no normal destination and includes init bytecode.
\item Execution outcome can modify balances, storage, and logs.
\item Failed execution reverts state changes but still spends gas.
\item Receipts record status, gas used, and emitted events.
\end{itemize}

\subsubsection*{Q2. Explain various elements present in Ethereum blockchain.}
\begin{itemize}
\item Ethereum block contains header and body.
\item Header includes parent hash, state root, transaction root, receipts root, timestamp, gas fields.
\item Transaction list contains user and contract interaction requests.
\item Receipt list records execution results for each transaction.
\item State root points to world state snapshot after block execution.
\item Transaction root binds integrity of all transactions in that block.
\item Receipts root binds integrity of all transaction results.
\item Logs and events provide structured outputs from contract execution.
\item Accounts include EOAs and contract accounts with different control models.
\item Contract bytecode and storage represent on-chain application backend.
\end{itemize}

\subsubsection*{Q3. Compare Bitcoin and Ethereum. Explain Ethereum world state with neat diagram.}
\begin{itemize}
\item Bitcoin primarily focuses on secure peer-to-peer digital money.
\item Ethereum focuses on programmable contracts plus value transfer.
\item Bitcoin scripting is intentionally limited; Ethereum supports richer logic.
\item Bitcoin commonly uses UTXO accounting model.
\item Ethereum commonly uses account-state model.
\item Bitcoin block interval is around 10 minutes; Ethereum has faster block cadence.
\item Ethereum world state is global mapping of all account states at block height.
\item Each processed block updates account balances, nonces, and contract storage.
\item Updated world state is represented by a new state root hash.
\item That root in block header proves state consistency for that block.
\item[]
\vspace{0.25em}
\noindent\framebox[0.97\linewidth][c]{
\begin{minipage}{0.93\linewidth}
\vspace{0.08cm}
\centerline{\textbf{World State Diagram}}
\centerline{Pre-State + Transactions}
\centerline{$\Downarrow$}
\centerline{EVM Execution}
\centerline{$\Downarrow$}
\centerline{Post-State + New State Root in Header}
\vspace{0.08cm}
\end{minipage}}
\end{itemize}

\subsubsection*{Q4. Explain how transactions are processed in Ethereum. Explain gas and its effect on processing.}
\begin{itemize}
\item User signs transaction in wallet and sends it to network.
\item Nodes perform basic checks and place valid transaction in mempool.
\item Validator/miner selects transactions, usually prioritizing fee incentives.
\item Transactions are executed by EVM in block order.
\item Contract calls may read/write storage and emit logs.
\item Execution consumes gas for each opcode and storage operation.
\item If gas limit is insufficient, execution reverts.
\item Reverted transaction still pays for computation already consumed.
\item Gas protects network from infinite loops and spam.
\item Higher offered fee generally improves inclusion speed.
\item[]
\vspace{0.25em}
\noindent\framebox[0.97\linewidth][c]{
\begin{minipage}{0.93\linewidth}
\vspace{0.08cm}
\centerline{\textbf{Processing Flow}}
\centerline{Sign $\rightarrow$ Broadcast $\rightarrow$ Mempool $\rightarrow$ Block Inclusion}
\centerline{$\Downarrow$}
\centerline{EVM Execute $\rightarrow$ Receipt + State Update}
\vspace{0.08cm}
\end{minipage}}
\end{itemize}

\subsubsection*{Q5. Explain control structures in Solidity.}
\begin{itemize}
\item Solidity supports conditional branches using \texttt{if}, \texttt{else if}, and \texttt{else}.
\item Loops include \texttt{for}, \texttt{while}, and \texttt{do-while}.
\item Loop usage should be controlled because gas usage grows with iterations.
\item \texttt{break} exits loop early and \texttt{continue} skips one iteration.
\item \texttt{return} exits function and can provide output value.
\item \texttt{require} validates inputs and reverts with message on failure.
\item \texttt{revert} manually rolls back state with optional reason.
\item \texttt{assert} is used for invariants and internal consistency checks.
\end{itemize}

\subsubsection*{Q6. Write Solidity bank contract (deposit, withdraw, view balance).}
\begin{itemize}
\item This contract maps each user address to its stored balance.
\item Deposit function accepts ETH and increases sender balance.
\item Withdraw function checks balance and transfers requested amount.
\item Balance view function returns sender's current stored value.
\item Guard checks prevent invalid withdraw attempts.
\item State update before transfer reduces reentrancy risk.
\end{itemize}

\begin{lstlisting}
// SPDX-License-Identifier: MIT
pragma solidity ^0.8.20;

contract SimpleBank {
    mapping(address => uint256) private balances;

    function deposit() external payable {
        require(msg.value > 0, "amount=0");
        balances[msg.sender] += msg.value;
    }

    function withdraw(uint256 amount) external {
        require(amount > 0, "amount=0");
        require(balances[msg.sender] >= amount, "insufficient");

        balances[msg.sender] -= amount;
        (bool ok, ) = msg.sender.call{value: amount}("");
        require(ok, "transfer failed");
    }

    function viewBalance() external view returns (uint256) {
        return balances[msg.sender];
    }
}
\end{lstlisting}

\subsubsection*{Q7. Define smart contract. Write Solidity contract to store and retrieve a single string.}
\begin{itemize}
\item A smart contract is code on blockchain that executes predefined rules.
\item It can store data and expose functions to update or read it.
\item For single string storage, one state variable is enough.
\item Setter function updates value.
\item Getter function returns current value.
\item Access control can be added if only owner should update.
\end{itemize}

\begin{lstlisting}
// SPDX-License-Identifier: MIT
pragma solidity ^0.8.20;

contract StringStore {
    string private value;

    function setValue(string calldata newValue) external {
        value = newValue;
    }

    function getValue() external view returns (string memory) {
        return value;
    }
}
\end{lstlisting}

\subsubsection*{Q8. Write Solidity voting contract.}
\begin{itemize}
\item A simple voting contract keeps vote count candidate-wise.
\item It tracks whether an address has already voted.
\item One-vote-per-address rule prevents duplicate voting.
\item Constructor can pre-load candidate list.
\item Vote function validates candidate id and voting status.
\item Result function returns current votes per candidate.
\end{itemize}

\begin{lstlisting}
// SPDX-License-Identifier: MIT
pragma solidity ^0.8.20;

contract SimpleVoting {
    string[] public candidates;
    mapping(uint256 => uint256) public votes;
    mapping(address => bool) public hasVoted;

    constructor(string[] memory names) {
        candidates = names;
    }

    function vote(uint256 candidateId) external {
        require(!hasVoted[msg.sender], "already voted");
        require(candidateId < candidates.length, "invalid candidate");

        hasVoted[msg.sender] = true;
        votes[candidateId] += 1;
    }

    function getVotes(uint256 candidateId) external view returns (uint256) {
        return votes[candidateId];
    }
}
\end{lstlisting}

\end{document}
